% % Chapter 1.1
% \selectlanguage{greek}
% \chapter{Περίληψη στα Ελληνικά} % Main chapter title

% \label{ChapterGreek}

% Το θέμα αυτής της εργασίας ήταν οι μέθοδοι αυτόματης περίληψης κειμένου με τη βοήθεια νευρωνικών δικτύων και πώς αυτές μπορούν να εφαρμοστούν σε επιστημονικά συγγράμματα. Η ανάπτυξη αξιόπιστων μεθόδων αυτόματης περίληψης είναι ένα πρόβλημα που τα τελευταία χρόνια έχει μελετηθεί εκτενώς στην βιβλιογραφία. Στο πρόβλημα της αυτόματης περίληψης κειμένων επιδιώκουμε να κατασκευάσουμε με αυτόματο τρόπο ένα σύντομο κείμενο, πολύ συντομότερο από το πρωτότυπο που όμως να περιλαμβάνει με ακρίβεια τα βασικά στοιχεία που περιγράφονται στο αρχικό κείμενο. Το πρόβλημα αυτό είναι σημαντικό απο πολλές απόψεις. Πρώτον, η αναπαράσταση ενός συγγράμματος με μια σύντομη περίληψη μπορεί να βελτιώσει σημαντικά την επίδοση μηχανών αναζήτησης και συστημάτων ανάκλησης πληροφορίας. Δεύτερον, μπορούν να διευκολύνουν έναν αναγνώστη να βρεί εύκολα ποιά συγγράμματα είναι σχετικά με τα ενδιαφέροντά του και ποιά όχι. Τέλος, θα ήταν εξαιρετικά χρήσιμο για έναν συγγραφέα που μόλις ολοκλήρωσε το άρθρο του να μπορεί να δεί μια αυτόματη περίληψη που στη συνέχεια μπορεί να τροποποιηθεί ώστε να δημοσιευτεί σαν περίληψη του άρθρου.

% Έχει αναπτυχθεί ένας αριθμός από μεθόδους που βασίζονται στα νευρωνικά δίκτυα, οι οποίες πετυχαίνουν καλά αποτελέσματα ειδικά όταν πρόκειται για σύντομα κείμενα. H πλειοψηφία αυτών των μεθόδων έχει εστιάσει σε σύντομα κείμανα, κυρίως άρθρα εφημερίδων και προσπαθούν να κατασκευάσουν μια περίληψη μερικών προτάσεων. Αντιθέτως, τα επιστημονικά συγγράμματα είναι αρκετά εκτενή και η περίληψη τους είναι αδύνατον να γίνει μέσα σε μερικές προτάσεις. Είναι κατανοητό ότι το συγκεκριμένο πρόβλημα είναι αρκετά πιο δύσκολο από την περίληψη άρθρων από εφημερίδες. Σε αυτή την εργασία μελετήσαμε πώς διάφορες σύγχρονες μέθοδοι πετυχαίνουν καλά αποτελέσματα σε σύντομα άρθρα θα μπορούσαν να εφαρμοστούν σε επιστημονικά συγγράμματα. 
% Αρχικά, στο κεφάλαιο \ref{neural_nets} επεξηγούμε τις βασικές αρχές των νευρωνικών δικτύων με έμφαση στα επαναλαμβανόμενο νευρωνικά δίκτυα που χρησιμοποιούνται κατα κόρον στη μοντελοποίηση κειμένων. Στη συνέχεια, στο κεφάλαιο \ref{neural_net_summarization} επεξηγούμε τη βασική αρχιτεκτονική κωδικοποιητή-αποκωδικοποιητή που χρησιμοποιήτε σε αρκετές από τις πιο σύγχρονες μεθόδους αυτόματης περίληψης και μηχανικής μετάφρασης. Τέλος, αναλύουμε τρείς βασικές μεθόδους που βασίζονται στο μοντέλο κωδικοποιητή-αποκωδικοποιητή για αυτόματη περίληψη κειμένων. Αυτές οι μέθοδοι στη συνέχεια θα εξαταστούν εκτενώς πειραματικά και θα μελετηθεί κατα πόσο μπορούν να κάνουν επιτυχώς αυτόματη περίληψη σε επιστημονικά συγγράμματα.

% Δημιουργήσαμε ένα νέο σύνολο δεδομένων με πάνω από 1 εκατομμύριο από επιστημονικά συγγράμματα στον τομέα της βιοιατρικής και των επιστημών υγείας. Για το κάθε σύγγραμα περιλαμβάνει το πλήρες κείμενο καθώς και την περίληψη του κειμένου γραμμένη από τους συγγραφείς του συγγράμματος. Η συλλογή των δεδομένων έγινε από το \foreignlanguage{english}{PubMed Central (PMC)} και την συλλογή των \foreignlanguage{english}{Open Access (OA)} άρθρων. Επίσης, δημιουργήσαμε ένα μικρότερο σύνολο δεδομένων με πάνω από 700 χιλιάδες συγγράματα που περιλαμβάνει το υποσύνολο των δεδομένων του \foreignlanguage{english}{PMC-OA} που έχουν περιλήψεις δομημένες  ενότητες. Αυτό το σύνολο μας επιτρέπει να αντιστοιχίσουμε κάθε ενότητα το κειμένου με την αντίστοιχη ενότητα της περίληψης και να εκπαιδεύσουμε μοντέλα μηχανικής μάθησης που κάνουν την περίληψη της κάθε ενότητας. Η συλλογή και επεξεργασία των δεδομένων ώστε να δημιουργηθούν τα δύο σύνολα, περιγράφεται εκτενώς στο κεφάλαιο \ref{datasets}.

% Όπως προαναφέρθηκε, υπάρχουν ορισμένες διαφορές μεταξύ των άρθρων απο εφημερίδες και των επιστημονικών συγγραμμάτων. Η πιο προφανής, είναι η διαφορά στο μέγεθος του κειμένου. Ένα άρθρο εφημερίδας συνήθως έχει έκταση μέχρι μία σελίδα, ενώ στην περίπτωση των επιστημονικών συγγραμμάτων εκτείνονται σε αρκετές σελίδες. Πέρα απο αυτό, τα επιστημονικά κείμενα περιλλαμβάνουν έναν αριθμό από σημαντικές πληροφορίες, όπως μέθοδοι πειραματικά αποτελέσματα και βιβλιογραφία, ενώ ένα άρθρο εφημερίδας είναι συνήθως εστιασμένο σε ένα βασικό γεγονός. Δεδομένων αυτών των διαφορών, είναι επόμενο η περίληψη επιστημονικών συγγραμμάτων να διαφοροποιείται και αυτή σε κάποιο βαθμό. Οι υπάρχουσες μέθοδοι που δουλεύουν καλά σε σύντομα κείμενα, βρήκαμε ότι αντιμετωπίζουν δυσκολίες όταν εφαρμόζονται ως έχουν σε σημαντικά μεγαλύτερα και ποιό πολύπλοκα κείμενα. Προκειμένου να αντιμετωπίσουμε αυτές τις δυσκολίες, αναπτύξαμε μια νέα μέθοδο εκπαίδευσης για τα μοντέλα που κάνουν αυτόματη περίληψη και την εφαρμόσαμε σε συνδυασμό με ένα αριθμό από σύγχρονες μεθόδους αυτόματης περίληψης στα δεδομένα που δημιουργήσαμε.

% Η μέθοδος αυτή, εκμεταλλεύεται το σύνολο δεδομένων με τις δομημένες περιλήψεις και δημιουργεί παραδείγματα εκπαίδευσης που περιλαμβάνουν μία ενότητα του κειμένου μαζί με την αντίστοιχη ενότητα της περίληψης. Αυτά τα παραδείγματα χρησιμοποιούνται με τέτοιο τρόπο ώστε να εκπαιδεύσουμε μοντέλα μηχανικής μάθησης ώστε να κάνουν περίληψη της κάθε ενότητας ξεχωριστά. Μόλις τα μοντέλα εκπαιδευτούν επαρκώς, μπορούν να πάρουν σαν είσοδο την κάθε ενότητα ενός συγγράμματος και να κατασκευάσουν μια περίληψη της ενότητας. Στη συνέχεια, όλες οι επιμέρους περιλήψεις μπορούν να συνδυαστούν ώστε να δημιουργήσουν την τελική περίληψη του άρθρου. Οι περίληψη που θα προκείψει, είναι προφανώς δομημένη σε ενότητες που περιλαμβάνουν τα βασικά στοιχεία της κάθε ενότητας. Με αυτό τον τρόπο, αντιμετωπίζουμε εν μέρη τη δυσκολία του προβλήματος, κατακερματίζοντας το σύγγραμα σε επιμέρους κείμενα που είναι σχετικά μικρά σε μέγεθος και αρκετά εστιασμένα.

% Κάναμε εκτενή πειράματα, όπου χρησιμοποιήσαμε τα τρία μοντέλα που περιγράφονται στο κεφάλαιο \ref{neural_net_summarization} στο καινούργιο σύνολο δεδομένων με τις δομημένες περιλήψεις. Δοκιμάσαμε να εφαρμόσουμε τις τρείς μεθόδους με δύο διαφορετικούς τρόπους. Αρχικά, επιχειρήσαμε να τις εφαρμόσουμε ως έχουν χρησιμοποιόντας το πλήρες κείμενο και προσπαθόντας να παράξουμε μια συνολική περίληψη. Αυτή η μέθοδος είναι αρκετά απλοϊκή και βρήκαμε ότι και οι τρείς μέθοδοι βρήκαν αρκετές δυσκολίες. Για να διευκολύνουμε τα μοντέλα επιλέξαμε να συμπεριλάβουμε στο κείμενο προς περίληψη μόνο τις ενότητες της εισαγωγής, περιγραφής των μεθόδων και των συμερασμάτων. Επίσης, περιορίσαμε το κείμενο εισόδου σε 500 λέξεις και το κείμενο εξόδου σε 120 λέξεις ώστε να είναι δυνατή η εκπαίδευση των μοντέλων. Εκτός από την απλοϊκή μέθοδο, δοκιμάσαμε να εφαρμόσουμε την μέθοδο με τις δομημένες περιλήψεις που περιγράψαμε προηγουμένως. Αυτή η μέθοδος δεν έχει τις δυσκολίες της απλοικής μεθόδου καθώς λειτουργεί σε κάθε ενότητα χωριστά. Και πάλι επιλέξαμε να συμπεριλάβουμε στο κείμενο προς περίληψη μόνο τις ενότητες της εισαγωγής, περιγραφής των μεθόδων και των συμερασμάτων και περιορίσαμε το κείμενο εισόδου σε 500 λέξεις και το κείμενο εξόδου σε 120 λέξεις για την κάθε ενότητα.

% Μέσα από τα πειράματα που έγιναν, βρήκαμε ότι η νέα αυτή μέθοδος μπορεί να αυξήσει τις επιδόσεις των κορυφαίων μεθόδον στο εν λόγω πρόβλημα. Πιο συγκεκριμένα, όλα τα μοντέλα φαίνονται να δυσκολεύονται όταν η εκπαίδευσή τους γίνεται πάνω σε όλο το κείμενο του συγγράματος. Είμαστε υποχρεωμένοι σε κάθε περίπτωση να κόβουμε το κείμενο σε 500 λέξεις καθώς ήταν σχεδόν αδύνατο να εκπαιδεύσουμε μοντέλα με μεγαλύτερη είσοδο. Αυτό έχει ως αποτέλεσμα να χάνουμε αρκετή και σημαντική πληροφορία, κάτι που έχει σοβαρές συνέπειες στην περίληψη που προκύπτει. Επίσης, αναγκαστικά περιορίζουμε την έξοδο σε 100-120 λέξεις που σημαίνει ότι είναι πολύ δύσκολο για τα μοντέλα να εστιάσουν στα διαφορετικά κομμάτια του κειμένου. Αντίθετα, όταν η εκπαίδευση γίνεται σε κάθε ενότητα του κειμένου ξεχωριστά, περιοριζόμαστε σε 500 λέξεις εισόδου για κάθε ενότητα με αποτέλεσμα να διατηρούμε πολύ μεγαλύτερο ποσοστό της πληροφορίας. Επίσης, επειδή η περίληψη πλέον γίνεται σε επίπεδο ενότητας, τα μοντέλα μπορούν εύκολα να εστιάσουν σε διαφορετικά κομμάτια του κειμένου. Ως αποτέλεσμα, οι περιλήψεις που προκύπτουν με αυτή τη μέθοδου είναι πολύ ανώτερης ποιότητας σε σχέση με την πιο απλοϊκή μέθοδο. Η πειραματική διαδικασία και τα αποτελέσματα της παρουσιάζονται και σχολιάζονται εκτενώς στο κεφάλαιο \ref{experiments}. Επίσης, συμπεριλαμβάνουμε και παραδείγματα που επιδεικνύουν την ικανότητα των διαφόρων μεθόδων να κατασκευάσουν καλής ποιότητας περιλήψεις. Τέλος, στο κεφάλαιο \ref{conclusion} σχολιάζουμε τα ευρήματα αυτής της εργασίας, και παραθέτουμε κατευθύνσεις για μελλοντική έρευνα.

\selectlanguage{english}
